\chapter{Projekt układu elektronicznego}\label{ch:elektronika}
% Ok. 18-22 stron. Rytm każdej sekcji: wymagania bloku -> wycinek schematu -> dobór elementów
% -> obliczenia ze wzorem i referencją do datasheetu -> co weryfikuje rozdz. 7.
\begin{zawartosc}
  \item Zasilanie: nPM1300 (BUCK2 3,3 V, ładowarka CC/CV, LOADSW1, tryb Ship), ochrona przed odwrotną polaryzacją; ogniwo bez NTC (R13 0 Ω zamiast termistora) i konsekwencje tej decyzji.
  \item Mikrokontroler nRF52840: konfiguracja zasilania, odsprzęganie, kwarce z obliczeniem pojemności obciążenia, programowanie SWD, przydział GPIO.
  \item Tor radiowy: dopasowanie anteny, linia 50 Ω (CPWG) na wybranym stackupie, antena w obudowie ze stalową płytą.
  \item Interfejsy: USB-C z zabezpieczeniami, łącze MB-DB, I2C, złącze rozszerzeń J3, dioda statusu.
  \item Membrana FPC: matryca 6 x 17 col2row, folia ROW (2 warstwy, z diodami) i folia COL (1 warstwa) na oryginalnym spacerze, złącza J4/J7.
  \item Projekt PCB (4 warstwy, rozmieszczenie, pętle prądowe, punkty pomiarowe) oraz BOM, koszt i produkcja w JLC (controller zamówiony 16.11, membrana 23.11).
\end{zawartosc}

\section{PMIC nPM1300}\label{sec:pmic}
% - BUCK2 = 3,3 V (VSET2 470k), BUCK1 off, VDDIO, I2C, GPIO0 przerwanie, GPIO1 LOADSW1.
% - Ładowarka CC/CV, bez NTC (R13 0 Ω, firmware thermistor-ohms = 0), tryb Ship, ochrona przed odwrotną polaryzacją.
% - Obliczenia: HD 'Blok PMIC: obliczenia' (tętnienia cewki, zakres regulacji, straty ładowarki, spadek na Q1).
% - Tabela połączeń: HD 'Blok PMIC: połączenia'.
\todo{Napisać.}

\section{MCU nRF52840}\label{sec:mcu}
% - Konfiguracja 5 (VDD = VDDH, REG1 DC/DC), odsprzęganie, HFXO/LFXO, SWD TC2030, RESET, SWO.
% - Obliczenia: HD 'Obliczenia: obciążenie kwarców' (historia błędnego CL, CL M1).
% - Przydział GPIO: HD 'Przydział GPIO', CL M7.
\todo{Napisać.}

\section{Tor RF}\label{sec:rf}
% - Dopasowanie 4-elementowe z config 7, U.FL, antena w obudowie ze stalową płytą.
% - CPWG 50 Ohm na JLC04121H-3313: 0,135/0,15 mm (CL A30, HD 'PCB: stackup'); L1 Murata (CL A19).
\todo{Napisać.}

\section{Interfejsy}\label{sec:interfejsy}
% - USB-C bez 5,1k, ESD, koordynacja zabezpieczeń VBUS (HD).
% - Łącze MB-DB 1x10 i pinout; I2C pull-upy (HD 'Obliczenie: pull-upy I2C').
% - J3 (HD 'Smart features'), dioda WS2812B przez LOADSW1 (HD 'Dioda statusu').
\todo{Napisać.}

\section{Matryca klawiszy i membrana FPC}\label{sec:membrana}
% - Matryca 6x17, col2row, dioda na klawisz, J4 (18) i J7 (6) (HD 'Złącza matrycy').
% - Folia ROW 2L z diodami + folia COL 1L na oryginalnym spacerze (hardware/membrane, zamówione 23.11, production/membrane).
\todo{Napisać.}

\section{Projekt PCB}\label{sec:pcb}
% - Stackup 4L 1,2 mm (CL A24), rozmieszczenie, pętla BUCK2 i NT1 (CL A22), via-in-pad POFV (A18, A25).
% - Reguły DRC (A29), TP i R3 (A15, A28), DB wyłamywany. Widoki warstw, render 3D.
\todo{Napisać.}

\section{BOM, koszt i produkcja}\label{sec:bom}
% - JLC Basic vs Extended (HD 'BOM pod montaż JLC'), części lutowane ręcznie (U2, J3, J5).
% - Zamówienie rev A (CL A31), tabela kosztów.
\todo{Napisać.}

